\section{Introduction}
\label{sec:introduction}

Public procurement records describe how public buyers solicit work, select suppliers, and
publish awards and contracts.  The Open Contracting Data Standard (OCDS) provides a common
release and record model for these disclosures, including parties, tenders, awards, contracts,
and linked documents \citep{ocds-standard}.  The resulting data can support questions such as
which suppliers serve multiple buyers, whether disclosure is complete for a class of notices,
or how activity changes across years.  These are not, however, ordinary document lookup
questions.  Their answers may depend on the role of an organisation in one record, a set of
organisations derived from that record, a second retrieval bound by the derived set, and a final
set, count, comparison, or ranking operation.

This distinction matters because retrieving \emph{relevant} evidence is not equivalent to
retrieving the \emph{complete} denotation required by an aggregate.  Let $R(q)$ be the complete
set of contract records satisfying question $q$, and let a top-$k$ retriever return
$\widehat{R}(q)\subset R(q)$.  Even when $\widehat{R}(q)$ contains convincing examples, it does
not in general follow that
\begin{equation}
 |\widehat{R}(q)|=|R(q)|, \qquad
 \sum_{r\in\widehat{R}(q)}v(r)=\sum_{r\in R(q)}v(r),
 \label{eq:retrieval-not-aggregation}
\end{equation}
or that a ranking computed from the retrieved subset has the same maximum.  Retrieval-centred
KGQA systems are effective at finding answer-bearing subgraphs \citep{sun-etal-2019-pullnet,
jiang2023unikgqa,luo2024rog}, but exact counts and compositional set operations additionally
need explicit scope, typed intermediate values, and exhaustive execution within the declared
data universe.

Direct language-model generation introduces a different failure mode.  A model can produce a
well-formed query that reverses buyer and supplier roles, projects a contract identifier when
the question requests an organisation name, or aggregates a non-additive monetary field.  Such
answers can remain fluent and superficially plausible.  Prior work therefore represents complex
questions with executable programs \citep{cao-etal-2022-kqa,wolfson-etal-2020-break}, constrains
formal generation \citep{scholak-etal-2021-picard,shu-etal-2022-tiara}, or uses execution and
environment feedback \citep{chen-etal-2021-retrack,huang-etal-2024-queryagent}.  These ideas
motivate our central design choice: the language model proposes a typed computation, while the
data system---not the model's prose---has final authority over factual values.

We implement this idea at two connected levels.  First, a compact model is trained to translate
procurement questions into a closed algebra containing record filters, distinct projections,
semi- and anti-joins, set operations, grouping, aggregation, comparison, and ranking.  Every
program is type checked and executed deterministically over a frozen record universe.  Second,
the deployment runtime adds a structured understanding stage and a conditional planner.  A
valid intent is compiled deterministically when possible; otherwise a graph planner produces a
typed variable-dependency plan.  The plan is grounded against allowed schema fields, executed
exhaustively, and passed through preflight, postflight, answer-sanity, and evidence checks.
Structured defects may trigger bounded replanning, but a clean empty result is terminal so that
the system cannot search repeatedly until it finds a preferred answer.

\begin{figure}[t]
  \centering
  \figureplaceholder{2.05in}{\textbf{Author-drawn Figure 1: visual thesis and running example.}\\
  Show one real procurement bridge question, the structured intent produced by the
  understander, the actually taken compiler/planner branch, the source-record set, emitted
  bridge values, bound target-record set, deterministic result, and release checks.  Every
  displayed value must come from one saved end-to-end trace; unused branches are grey.}
  \caption{Intended content for the introduction figure.  The final artwork will instantiate a
  single saved end-to-end trace rather than combine artifacts from the separate compositional
  evaluation protocol.  Runtime verification denotes structural, execution, and evidence
  checks; oracle equality is an offline evaluation criterion.}
  \label{fig:visual-thesis}
\end{figure}

The distinction between the two levels is methodological rather than cosmetic.  Our primary
compositional evaluation asks whether a single model call emits a valid algebra tree or an
appropriate abstention, after which a deterministic evaluator computes the answer and an
offline scorer compares it with a frozen oracle.  It does not invoke the deployment
understander, graph planner, evidence bundle, or repair loop.  A separate system evaluation
measures the complete runtime.  Keeping these protocols separate prevents an offline oracle
check from being mislabeled as an inference-time verifier and prevents a component that was not
called for a test item from receiving credit for that item.

Experiments use a provenance-preserving graph constructed from UK procurement releases from
2022--2026 and a frozen \PacsN{}-question compositional test.  The typed-program model attains
\PacsOursA\% strict accuracy on the canonical question surface, a gain of \PacsGainBase{}
percentage points over the same untuned 8B base model and \PacsGainTeacher{} points over a cloud
planner.  Accuracy is \PacsOursB\% after independent surface naturalisation, a drop of only
\PacsSurfaceDrop{} points.  The strongest gains occur on overlap/exclusion and relational
composition, whereas supplier-concentration and disclosure/compliance remain substantially
harder.  A separate balanced \RuntimeN{}-item artifact evaluates the complete runtime and records
\RuntimeOurs\% accuracy, compared with \RuntimeBase\% for the untuned system.  We treat the
latter as secondary evidence because its stored outputs are less portable than the PACS
artifacts.

The contributions of this work are therefore:
\begin{itemize}
  \item a provenance-preserving procurement graph and a deterministic query-record projection,
  with explicit entity-resolution, additivity, and coverage checks;
  \item a closed typed algebra and program-first supervision pipeline for exact set, join,
  aggregation, comparison, and ranking questions;
  \item a full question-answering runtime that separates learned understanding and planning
  from grounding, exhaustive execution, evidence checks, bounded repair, and controlled
  non-answering; and
  \item PACS, a frozen compositional evaluation with two surface channels, operator-family and
  depth diagnostics, answerability controls, and an audit of the limitations of its isolation
  and scoring procedures.
\end{itemize}

The contribution is deliberately narrower than open-domain question answering.  The pipeline
does not guarantee that the source data are complete in the world, that stored organisation
names are unique legal entities, or that passing an execution check proves semantic correctness.
It instead makes a more testable claim: within a declared frozen data universe, the model's
proposed computation can be type checked, executed, traced, and rejected when its implemented
preconditions fail.
